\documentclass[11pt,a4paper]{article}
\usepackage[T1]{fontenc}
\usepackage{lmodern,microtype,amsmath,amssymb,amsthm,booktabs}
\usepackage{graphicx,xcolor,fancyhdr,float}
\definecolor{ink}{HTML}{193B50}
\usepackage{enumitem}
\setlist{itemsep=2pt,topsep=4pt}
\emergencystretch=1.5em
\setlength{\headheight}{14pt}
\pagestyle{fancy}\fancyhf{}
\fancyhead[L]{\small\itshape Moment thresholds for exact causal randomness}
\fancyhead[R]{\small Research manuscript}
\fancyfoot[C]{\small\thepage}
\renewcommand{\headrulewidth}{0.3pt}
\usepackage[margin=25mm]{geometry}
\usepackage[colorlinks=true,linkcolor=ink,urlcolor=ink,citecolor=ink]{hyperref}
\newtheorem{theorem}{Theorem}
\newtheorem{lemma}[theorem]{Lemma}
\newtheorem{corollary}[theorem]{Corollary}
\newtheorem{proposition}[theorem]{Proposition}
\newtheorem{conjecture}[theorem]{Conjecture}
\numberwithin{equation}{section}
\newcommand{\E}{\mathbb E}
\newcommand{\Prb}{\mathbb P}
\newcommand{\Law}{\operatorname{Law}}
\newcommand{\TV}{\mathrm{TV}}
\newcommand{\ii}{\imath}
\title{\bfseries Moment thresholds for exact causal randomness\\[6pt]
\large Information spectra, adaptive envelopes, and a logarithmic coefficient}
\author{}
\date{26 September 2026}
\begin{document}
\maketitle
\begin{abstract}
Consider exact simulation of two finite output laws, selected adaptively
for $n$ rounds, using a single initial discrete seed. For equal-entropy
pure dyadic laws, we prove that the minimum seed entropy is $nh$ plus
an excess of order $\sqrt n$ when information variances differ, of order
$\log n$ when only the third centered information moments first differ,
and bounded when both moments agree. The result follows from exact
attainment of a recursive information-spectrum bound and a moment
classification of controlled random-walk envelopes. The latter
classification holds for arbitrary finite centered increment laws.
The proof uses an exact $L^1$ entropy-increment identity, a Fourier lower
bound, and an integrated envelope contraction with compact smoothing.
We also formulate a sharp logarithmic-coefficient conjecture, isolate a
sufficient $L^1$ remainder estimate, and report reproducible computations
for seven dyadic families. The coefficient remains unproved.
\end{abstract}
\noindent\textbf{Keywords.} Exact simulation; minimum-entropy coupling;
information spectrum; controlled random walk; moment matching.

\section{Introduction}
The entropy needed to supply independent outputs of one prescribed law
is additive. A different question arises when the law can be selected
after observing earlier outputs, while all randomness must be supplied
at time zero. Equal output entropies remove the first-order distinction
between the actions. They do not remove the cost of supporting both
choices at every observable history.

This paper studies that excess cost in a history-indexed model. Its main
result identifies a finite moment threshold: for pure dyadic action laws,
matching information variance and the third centered information moment
is necessary and sufficient for uniformly bounded excess. A third-moment
defect produces a logarithm; a variance defect produces a square root.
The theorem classifies every finite pure dyadic pair of equal entropy,
without symmetry assumptions. It does not provide a uniform constant as
the pair varies.

The proof separates two questions. First, a recursive spectral bound is
exactly realizable by a seed in the pure dyadic class. Second, the mean
of that spectral envelope can be classified analytically. The analytic
part does not require a lattice or dyadic probabilities. This distinction
also identifies the obstruction to extending the causal upper bounds:
for general action laws, an envelope need not be the information spectrum
of an exactly realizable seed.

\paragraph{Relation to existing work.}
Minimum-entropy coupling and functional representation provide the
one-round setting. Majorization-based constructions and approximation
bounds were developed by Cicalese, Gargano, and Vaccaro~\cite{CGV19}
and Li~\cite{Li21}. Information-spectrum converses for coupling and
functional representation were established by Shkel and
Yadav~\cite{SY23}; Compton et al.~\cite{CKQGK23} developed a profile
lower bound beyond the majorization bound. The spectral ordering used
below belongs to this established framework. Our argument applies it
recursively under conditional, history-by-history constraints.

The approximation landscape has also changed recently. Compton's
additive-$\varepsilon$ scheme~\cite{Compton26} is polynomial in support
size for a fixed number of marginals and fixed $\varepsilon$.
Minimum-R\'enyi entropy coupling~\cite{Ma25} and the geometry of R\'enyi
entropy on the majorization lattice~\cite{YS26} address related static
questions. These results do not, as stated, give a horizon-uniform error
bound for the repeated conditional construction considered here. The
claims proved in this paper concern exact recursive attainment and the
asymptotics of its adaptive envelope; they are not claims to originate
majorization or information-spectrum lower bounds.

\paragraph{Martingale and nonlinear limit theory.}
There is a second relevant literature. The controlled sums below are
martingales, and logarithmic losses in martingale normal approximation
are well established; Guo~\cite{Guo24} gives recent Wasserstein bounds
and a detailed account of earlier results, including Bolthausen's sharp
Kolmogorov rate. Fang, Peng, Shao, and Song~\cite{FPSS19} obtain
$(1+\log n)/\sqrt n$ bounds for Lipschitz tests under sublinear
expectations. Thus a logarithm arising from adaptive increments is not
by itself an unfamiliar phenomenon. Our quantity is the integral of a
pointwise envelope over policies. The optimizing policy can vary with
the integration threshold, so a bound on the Wasserstein distance of
each individual controlled law does not directly bound this quantity.
The integrated contraction below deals with that order of optimization
and integration explicitly.

\paragraph{Organization and status.}
Sections~\ref{sec:main}--\ref{sec:closing} give the classification and its
proof. Section~\ref{sec:explicit} gives an explicit logarithmic pair with
all-horizon numerical constants. Section~\ref{sec:coefficient} separates
a conjectured coefficient from a proved conditional criterion, and
Section~\ref{sec:numerics} reports finite computations. The manuscript
makes no claim of established literature priority. In particular, the
sharp coefficient is not included in the classification theorem.

\section{Model and main results}\label{sec:main}

All information quantities use base-two logarithms. For a finite law $A$,
write $\ii_A(x)=-\log_2 A(x)$ on its positive support. Define
\[
h_A=\E\ii_A(X),\quad
v_A=\E(\ii_A(X)-h_A)^2,\quad
\kappa_A=\E(\ii_A(X)-h_A)^3,\qquad X\sim A.
\]
A \emph{pure dyadic law} has each positive atom equal to $2^{-k}$ for an
integer $k\ge0$. The restriction is to individual powers of two, not to
arbitrary rational numbers with a power-of-two denominator.

A horizon-$n$ simulator consists of one discrete seed $S$, independent
of any policy randomization, and deterministic maps
\[
Y_t=f_t(S,A_1,Y_1,\ldots,A_{t-1},Y_{t-1},A_t),\qquad 1\le t\le n,
\]
where $A_t\in\{P,Q\}$. A policy selects $A_t$ using the observed
history (and, if desired, independent private randomness). Admissibility
means
\[
\Law(Y_t\mid A_1,Y_1,\ldots,A_{t-1},Y_{t-1},A_t)=A_t
\]
at every positive-probability conditioning event, for every such policy.
There is no fresh simulator randomness after time zero. It suffices to
check deterministic policies: conditioning on a policy's private
randomness gives the general case. The policy observes past actions and outputs, not the seed.
Let $C_n^{\mathrm{causal}}(P,Q)$ be the infimum seed entropy at horizon $n$.
The seed is allowed to depend on $n$.
No identification of responses at different histories is imposed merely
because the same action appears there. All positive atom probabilities
are included and zero-probability symbols are omitted. We use $\ln$ for
the natural logarithm, including in the sharp coefficient below.

\begin{theorem}[Complete equal-entropy classification for pure dyadic laws]
\label{thm:classification}
Let $P,Q$ be finite pure dyadic laws with $h_P=h_Q=h$, and set
$\Delta_n=C_n^{\mathrm{causal}}(P,Q)-nh$.
The optimum is attained by an explicit finite seed. The following cases
are exhaustive:
\[
\boxed{
\begin{array}{ll}
v_P\ne v_Q &: \Delta_n=\Theta(\sqrt n),\\[2pt]
v_P=v_Q,\ \kappa_P\ne\kappa_Q &: \Delta_n=\Theta(\log(n+1)),\\[2pt]
v_P=v_Q,\ \kappa_P=\kappa_Q &: \Delta_n=O(1).
\end{array}}
\]
In the bounded case, $\Delta_n$ is nondecreasing and converges to a finite
limit. This limit is positive unless $P,Q$ are permutations of one another;
in that exceptional case $\Delta_n=0$ for all $n$.
All asymptotic constants may depend on the fixed pair.
\end{theorem}
Thus, within this class, the second and third centered information
moments completely determine the growth order. Higher moments may affect
the bounded limit but do not create further unbounded orders.

\begin{corollary}[Necessary conditions without dyadic assumptions]
\label{cor:general}
For arbitrary finite action laws of equal entropy, unequal information
variances imply $\Delta_n=\Omega(\sqrt n)$. Equal variances and unequal
third centered information moments imply $\Delta_n=\Omega(\log(n+1))$.
In particular, bounded overhead requires matching both moments.
\end{corollary}

\begin{theorem}[Moment classification of adaptive envelopes]\label{thm:analytic}
Let $\alpha,\beta$ be any finite centered probability laws on $\mathbb R$.
Define
\[
G_0=\mathbf1_{[0,\infty)},\qquad
G_{n+1}=\min\{\alpha*G_n,\beta*G_n\},
\]
and let $e_n$ be the mean of the probability law with CDF $G_n$.
Then $e_n$ is nonnegative and nondecreasing, and
\[
e_n=\begin{cases}
\Theta(\sqrt n),& \int x^2\alpha(dx)\ne\int x^2\beta(dx),\\
\Theta(\log(n+1)),& \int x^2\alpha(dx)=\int x^2\beta(dx),\quad
                    \int x^3\alpha(dx)\ne\int x^3\beta(dx),\\
O(1),& \int x^j\alpha(dx)=\int x^j\beta(dx),\quad j=2,3.
\end{cases}
\]
In the last case, the limiting mean is positive if $\alpha\ne\beta$,
and zero otherwise. The theorem imposes no arithmetic restriction on
supports or probabilities.
\end{theorem}
For general finite equal-entropy action laws, this theorem classifies a
lower bound on causal seed entropy. Exact causal attainment is the
additional dyadic statement in Lemma~\ref{lem:optimal}.

\section{Exact profile recursion and dyadic realization}
For any discrete law $A$, let $F_A(t)=\Prb\{\ii_A(X)\le t\}$, and let
$\mu_A=\Law(\ii_A(X))$. Define CDFs recursively by
\begin{equation}\label{eq:F}
F_0(t)=\mathbf1_{\{t\ge0\}},\qquad
F_n(t)=\min\{(\mu_P*F_{n-1})(t),(\mu_Q*F_{n-1})(t)\}.
\end{equation}
Convolution means $(\mu*F)(t)=\int F(t-x)\,\mu(dx)$.
The minimum of two finite-support CDFs is again a finite-support CDF.

\begin{lemma}[Profile lower bound and exact dyadic attainment]\label{lem:optimal}
For any finite $P,Q$, every admissible discrete horizon-$n$ seed $S$
satisfies $F_S(t)\le F_n(t)$ for all real $t$.
For pure dyadic $P,Q$, $F_n$ is the information CDF of a constructible
finite pure dyadic seed $R_n$ and
$C_n^{\mathrm{causal}}=H(R_n)$.
\end{lemma}
\begin{proof}
At horizon zero every seed has nonnegative information. Suppose the
profile bound holds for $n-1$ and fix the first action $a$ with output $Y$.
For every $y$ of positive probability, $S\mid Y=y$ is an admissible seed
for the remaining horizon. On that fiber,
\[
-\log_2\Prb\{S=s\}=
\ii_a(y)-\log_2\Prb\{S=s\mid Y=y\}.
\]
Induction gives $F_S(t)\le(\mu_a*F_{n-1})(t)$ for both first actions,
proving the lower bound. Integrating survival functions bounds seed
entropy below by the mean of $F_n$.

For finite pure dyadic laws $R,A$, a deterministic map from $R$ to $A$
exists if and only if $F_R(k)\le F_A(k)$ for all integers $k$.
Necessity follows because a source atom must fit inside its target atom.
For sufficiency, assign source atoms in decreasing order of size. At
mass $2^{-k}$, all eligible remaining target capacities are integer
multiples of $2^{-k}$ and total $F_A(k)-F_R(k-1)$, enough for the current
level. Thus all bins can be filled exactly.

For $F(k)=\min(F_A(k),F_B(k))$, the counts
$N_k=2^k(F(k)-F(k-1))$ are nonnegative integers. Taking $N_k$ atoms of
mass $2^{-k}$ realizes $F$ and refines both $A,B$.
Apply this recursively to $P\otimes R_{n-1}$ and $Q\otimes R_{n-1}$,
starting at a point mass $R_0$. The refinement maps send the current
state to an output and a new state. Conditional on that output, the new
state has exactly law $R_{n-1}$. Induction proves correctness after every
history and for every adaptive policy. The profile lower bound proves
optimality, even against non-dyadic seeds.
\end{proof}

\section{The envelope and its exact entropy increments}
Assume henceforth that both entropies equal $h$. Put
$\alpha=\Law(\ii_P-h)$ and $\beta=\Law(\ii_Q-h)$.
Both laws are centered. Let a controller choose either increment law
predictably, and let $S_j^\pi$ be its sum after $j$ steps. Define
\[
G_j(x)=\inf_\pi\Prb\{S_j^\pi\le x\},\qquad
U_j(x)=\sup_\pi\Prb\{S_j^\pi\le x\}.
\]
Dynamic programming, conditioning on the first increment, gives
\[
G_j=\min(\alpha*G_{j-1},\beta*G_{j-1}),\qquad
G_0=\mathbf1_{[0,\infty)}.
\]
The optimizing policy may depend on the threshold $x$; the envelope
need not be the CDF of any single controlled walk.
Consequently $G_j(x)=F_j(jh+x)$.
Let $Z_j$ have CDF $G_j$ and put $e_j=\E Z_j$.
For all finite action laws $\Delta_j\ge e_j$; for pure dyadic laws,
Lemma~\ref{lem:optimal} gives equality.
The two CDFs in the minimum have the same mean $e_{j-1}$. Therefore
\begin{equation}\label{eq:increments}
e_{j+1}-e_j=\frac12\| (\alpha-\beta)*G_j\|_1\ge0.
\end{equation}
Indeed, use $\min(A,B)=(A+B-|A-B|)/2$ and the CDF formula for differences
of means. This identity charges a nonnegative cost at every horizon.

Write $B=\max\{|x|:x\in\operatorname{supp}\alpha\cup
\operatorname{supp}\beta\}$. If $B=0$, both information laws coincide
and the problem is trivial. Assume $B>0$. For every policy, conditional
centering and $|D_i|\le B$ imply
\[
\Prb\{S_j^\pi\ge t\},\ \Prb\{S_j^\pi\le-t\}
\le\exp\left(-\frac{t^2}{2jB^2}\right),\qquad t>0.
\]
For completeness, convexity gives
$\E[e^{sD_i}\mid\mathcal F_{i-1}]\le\cosh(sB)\le e^{s^2B^2/2}$;
iterate and optimize the exponential Markov bound.
The estimate is uniform over policies, hence also bounds both tails of
$Z_j$. Integrating the tails yields
\begin{equation}\label{eq:second}
\E Z_j^2\le4B^2j,\qquad 0\le e_j\le2B\sqrt j.
\end{equation}

\section{A Fourier lower bound: why a logarithm is unavoidable}
Let $D=\alpha-\beta$, and use the Fourier convention
$\widehat\nu(t)=\int e^{itx}\nu(dx)$.
Suppose $d\ge2$ is the first order at which the moments of $\alpha,\beta$
differ. Finite support and Taylor expansion imply that for some
$c_0,t_0>0$,
\begin{equation}\label{eq:fourierD}
|\widehat D(t)|\ge c_0|t|^d\qquad(0<|t|\le t_0).
\end{equation}
Choose $c=\min(t_0,1/(2B))$ and $t_j=c/\sqrt{j+1}$.
Equation~\eqref{eq:second} and $\cos u\ge1-u^2/2$ show
\[
|\widehat{\Law(Z_j)}(t_j)|\ge
\Re\widehat{\Law(Z_j)}(t_j)
\ge1-\tfrac12t_j^2\E Z_j^2\ge\tfrac12.
\]
The function $R_j=D*G_j$ is compactly supported, and its distributional
derivative is $D*\Law(Z_j)$. Integration by parts therefore gives
\[
\|R_j\|_1\ge|\widehat R_j(t_j)|
=\frac{|\widehat D(t_j)|\,
|\widehat{\Law(Z_j)}(t_j)|}{|t_j|}
\ge\frac{c_0c^{d-1}}2(j+1)^{-(d-1)/2}.
\]
Together with \eqref{eq:increments}, this proves
\begin{equation}\label{eq:lower}
e_n\ge\frac{c_0c^{d-1}}4
\sum_{j=0}^{n-1}(j+1)^{-(d-1)/2}.
\end{equation}
For $d=2$ this is $\Omega(\sqrt n)$; for $d=3$ it is
$\Omega(\log(n+1))$. Since $\Delta_n\ge e_n$ without a dyadic
assumption, Corollary~\ref{cor:general} follows already.
The upper estimate in \eqref{eq:second} proves the matching square-root
upper bound for the envelope, and hence for the causal excess in the
pure dyadic class.

\section{An integrated envelope contraction}
For the remaining upper bounds, put
\[
W_j=\int_{\mathbb R}(U_j-G_j)(x)\,dx.
\]
The all-$\beta$ walk has mean zero and its CDF lies between the envelopes.
Thus $0\le e_j\le W_j$. No reflection symmetry is assumed here.
To distinguish this width from an ordinary Wasserstein bound, for any
reference CDF $B_j$ lying between the envelopes one has
\[
\sup_\pi\int|F_j^\pi-B_j|\,dx
\;\le\;\int\sup_\pi|F_j^\pi-B_j|\,dx
\;\le\;W_j.
\]
In general the first inequality cannot be reversed. In particular,
a lower bound on $e_j$ does not establish the same lower bound on the
Wasserstein error of one controlled walk.

Let $V$ be independent smoothing noise supported in $[-A,A]$, let $T_r$
be the all-$\beta$ walk, and define
\[
h_r(x)=\Prb\{T_r+V\le x\},\quad
R_r=D*h_r,\quad K_r=|R_r|,\quad
 a_r=\|K_r\|_1,\quad b_r=\|K_r'\|_1.
\]
The noise chosen below ensures that these functions are absolutely
continuous and compactly supported where needed.

\begin{lemma}[Contraction]\label{lem:contraction}
For every $n\ge1$,
\begin{equation}\label{eq:contract}
W_n\le2A+2\sum_{r=0}^{n-1}(a_r+b_rW_{n-1-r}).
\end{equation}
If $b_*:=\sum_{r\ge0}b_r<1/2$, then for every $N\ge1$,
\begin{equation}\label{eq:maxW}
\max_{0\le n\le N}W_n
\le\frac{2A+2\sum_{r=0}^{N-1}a_r}{1-2b_*}.
\end{equation}
\end{lemma}
\begin{proof}
Write $G_n^V,U_n^V$ for the smoothed policy envelopes and $W_n^V$ for
their integrated width. Bounded support of $V$ gives
$U_n(x)\le U_n^V(x+A)$ and $G_n(x)\ge G_n^V(x-A)$.
Every CDF $F$ satisfies $\int(F(x+A)-F(x-A))dx=2A$ by Tonelli, so
$W_n\le W_n^V+2A$.

For each policy, telescope $\E h_{n-i}(x-S_i^\pi)$ over $i$.
Since $h_{r+1}=\beta*h_r$, conditioning on the past gives
\[
\Prb\{S_n^\pi+V\le x\}-h_n(x)
=\sum_{i=1}^n\E\left[
\mathbf1_{\{\text{action }i\text{ uses }\alpha\}}
R_{n-i}(x-S_{i-1}^\pi)\right].
\]
Its absolute value is at most
$\sum_i\E K_{n-i}(x-S_{i-1}^\pi)$.
Let $B_j$ be the CDF of $T_j$. For every controlled prefix CDF $F_j^\pi$,
$|F_j^\pi-B_j|\le U_j-G_j$. Stieltjes integration by parts gives
\[
\E K_r(x-S_j^\pi)\le\E K_r(x-T_j)
+\int|K_r'(x-y)|(U_j(y)-G_j(y))\,dy.
\]
The right side is independent of $\pi$, so taking the supremum first
and then integrating in $x$ proves
\[
\int\sup_\pi\E K_r(x-S_j^\pi)dx\le a_r+b_rW_j.
\]
The smoothed envelope width is at most twice the largest absolute
deviation from $h_n$ at each threshold. Summing proves
\eqref{eq:contract}. Taking a maximum over $n\le N$, with $W_0=0$,
and moving $2b_*$ times that finite maximum to the left proves
\eqref{eq:maxW}.
\end{proof}

\section{Derivative estimates for an arbitrary finite reference law}
Assume $\beta$ is nondegenerate. Choose two of its support points
$x<y$, put $b=(y-x)/2>0$ and
$\rho=2\min(\beta\{x\},\beta\{y\})>0$.
Then $\beta=\rho\eta+(1-\rho)\zeta$, where
$\eta=(\delta_x+\delta_y)/2$ and $\zeta$ is a probability law when
$\rho<1$. The case $\rho=1$ needs no residual law.
Let $v_a$ denote the uniform density on $[-a,a]$, choose
$a=Lb$ with a positive integer $L$, and take $V$ to be the sum of
four independent variables with density $v_a$. Then $A=4a$ and, in
weak-derivative notation,
\begin{equation}\label{eq:derivative}
h_r^{(q+1)}=(v_a')^{*q}*v_a^{*(4-q)}*\beta^{*r},
\qquad q=2,3,4.
\end{equation}
Here $v_a'=(\delta_{-a}-\delta_a)/(2a)$, and exponent zero means
$\delta_0$. In particular $\|h_r^{(q+1)}\|_{\TV}\le a^{-q}$.
The fourth case can be a signed measure; for densities the total
variation norm agrees with the $L^1$ norm.

\begin{lemma}\label{lem:decay}
There are constants $C_q<\infty$ depending on $\beta,q$, but independent
of $a=Lb$ and $r$, such that
\begin{equation}\label{eq:decay}
\|h_r^{(q+1)}\|_{\TV}
\le\min\{a^{-q},C_q(r+1)^{-q/2}\},\qquad q=2,3,4.
\end{equation}
\end{lemma}
\begin{proof}
Telescoping point masses shows
$v_a'=\nu_L*v_b'$, where $\nu_L$ is uniform on
$\{-a+b,-a+3b,\ldots,a-b\}$.
Convolution with probability measures cannot increase total variation.
In the binomial mixture expansion of $\beta^{*r}$, the number $K$ of
$\eta$ increments has law $\operatorname{Bin}(r,\rho)$.
For a fixed $K=k$, translation and scaling of symmetric signs give
\[
\|(v_b')^{*q}*\eta^{*k}\|_{\TV}
\le b^{-q}\binom{k+q}{q}^{-1/2}.
\]
To verify this inequality, let $N=k+q$ independent symmetric signs be
$\varepsilon_i$. Up to sign and the factor $b^{-q}$, the signed measure
is the law of their sum weighted by $\prod_{i=1}^q\varepsilon_i$.
Given the sum, the conditional product equals the average of
$\prod_{i\in I}\varepsilon_i$ over all $q$-element subsets $I$.
These products are orthonormal, so that average has squared $L^2$ norm
$\binom Nq^{-1}$. Cauchy--Schwarz proves the asserted variation bound.

It remains to average over $K$. Using
$\binom{k+q}{q}\ge(k+1)^q/q!$ and splitting at $K=\rho r/2$ gives
\[
\E\binom{K+q}{q}^{-1/2}
\le\sqrt{q!}(1+\rho r/2)^{-q/2}+e^{-\rho r/8}
\le C'_q(r+1)^{-q/2}.
\]
For the exponential term, apply Markov to $e^{-(\ln2)K}$ and use
$1-\rho/2\le e^{-\rho/2}$ and $\ln2<3/4$.
This proves the decay estimate; combine it with the direct bound
$a^{-q}$ from \eqref{eq:derivative}.
\end{proof}

\section{Closing the classification}\label{sec:closing}
The remaining step is a cancellation estimate that does not use a
lattice. Suppose the moments of $D=\alpha-\beta$ vanish through order
$m-1$. Taylor's formula with integral remainder gives
\begin{equation}\label{eq:factor}
D*f=\Gamma_m*f^{(m)},\qquad
\Gamma_m=\int \frac{(-y)^m}{(m-1)!}
                 \int_0^1(1-t)^{m-1}\delta_{ty}\,dt\,D(dy).
\end{equation}
The terms of degrees below $m$ vanish after integration against $D$.
The signed measure $\Gamma_m$ is finite and compactly supported, with
\[
\|\Gamma_m\|_{\TV}\le M_m:=\frac1{m!}\int |y|^m\,|D|(dy).
\]
The identity first holds for smooth $f$; convolution with a smooth
approximate identity extends it to the weak derivatives used here.
For the smoothing of Lemma~\ref{lem:decay}, $R_r=D*h_r$ is absolutely
continuous and its derivative is integrable. Since
$\|(|R_r|)'\|_1\le\|R_r'\|_1$, we obtain
\begin{equation}\label{eq:ab}
a_r\le M_m\|h_r^{(m)}\|_{\TV},\qquad
b_r\le M_m\|h_r^{(m+1)}\|_{\TV}.
\end{equation}
If the common variance is positive, the reference law is nondegenerate and
Lemma~\ref{lem:decay} applies.
For $m=3$ or $4$, the summable sequence
$M_mC_m(r+1)^{-m/2}$ dominates the bounds on $b_r$.
As $a=Lb\to\infty$, their pointwise upper bound $M_ma^{-m}$ tends to
zero. Dominated convergence lets us fix a finite $a$ with
$\sum b_r<1/2$.

If variances match but third moments differ, use $m=3$.
Then $a_r\le M_3C_2/(r+1)$, and \eqref{eq:maxW} gives
$W_n=O(\log(n+1))$. Since $e_n\le W_n$, the Fourier lower bound proves the logarithmic case
of Theorem~\ref{thm:analytic}.

If variances and third moments both match, use $m=4$.
Then $a_r\le M_4C_3(r+1)^{-3/2}$ is summable, so
\eqref{eq:maxW} gives $W_n=O(1)$ and $e_n=O(1)$.
If the common variance is zero, both centered increment laws are $\delta_0$ and this conclusion holds trivially.
Finally \eqref{eq:increments} proves monotonicity. When
$\alpha\ne\beta$, it also gives
$e_1=\tfrac12\int|F_\alpha-F_\beta|>0$.
This completes Theorem~\ref{thm:analytic}.
For pure dyadic action laws, $\Delta_n=e_n$ by Lemma~\ref{lem:optimal}.
Equality of the information laws is equivalent to
equality of all atom multiplicities, hence to being permutations of one
another. In that case the recursion is the ordinary product profile and
$\Delta_n=0$. This completes Theorem~\ref{thm:classification}.

\section{An explicit logarithmic pair}\label{sec:explicit}
Take
\begin{align*}
P&=\bigl(1/16,(1/32)^{\times10},(1/64)^{\times12},
                       (1/128)^{\times56}\bigr),\\
Q&=\bigl((1/32)^{\times16},(1/128)^{\times64}\bigr).
\end{align*}
Their support sizes are $79$ and $80$. Their centered information laws are
\[
\alpha=\tfrac1{16}\delta_{-2}+\tfrac5{16}\delta_{-1}
       +\tfrac3{16}\delta_0+\tfrac7{16}\delta_1,
\qquad \beta=\tfrac12\delta_{-1}+\tfrac12\delta_1.
\]
Direct calculation gives
\[
h_P=h_Q=6,\qquad v_P=v_Q=1,\qquad
\kappa_P=-\tfrac38,\quad\kappa_Q=0,\qquad
D=\tfrac1{16}\delta_{-2}*(\delta_0-\delta_1)^{*3}.
\]
Let $\mathsf H_n=\sum_{k=1}^n1/k$ be the harmonic number.

\begin{theorem}[A genuine logarithmic overhead]\label{thm:explicit}
For the displayed fixed pair and every $n\ge1$,
\begin{equation}\label{eq:explicit}
\boxed{\frac{\mathsf H_n}{2048}
\le C_n^{\mathrm{causal}}(P,Q)-6n
<\frac{128}{5}+\frac3{10}\mathsf H_n.}
\end{equation}
In particular, the overhead is $\Theta(\log n)$ and is unbounded.
The exact one-step excess is $\Delta_1=1/8$.
\end{theorem}
\begin{proof}
For the lower bound, $B=2$ in \eqref{eq:second}. Choose
$t_j=1/(4\sqrt{j+1})$, so $|\widehat{\Law(Z_j)}(t_j)|\ge1/2$.
For $0\le t\le1/4$, the inequality
$2\sin(t/2)\ge t-t^3/24$ implies
$|1-e^{it}|^3\ge t^3/2$. Thus
$|\widehat D(t_j)|\ge t_j^3/32$, and
\eqref{eq:increments} gives
\[
e_{j+1}-e_j\ge\tfrac12\frac{t_j^3/32}{t_j}\cdot\tfrac12
=\frac1{2048(j+1)}.
\]
Summation proves the lower bound.

For the upper bound use four uniform smoothing variables on $[-2,2]$,
so $a=2$ and $A=8$. For the Rademacher reference, the derivative estimate
has the explicit form
\[
\|h_r^{(q+1)}\|_{\TV}
\le\min\{2^{-q},\binom{r+q}{q}^{-1/2}\}.
\]
Here the cubic finite-difference factorization of $D$ is sharper than
the general Taylor bound: the third unit difference is the third weak
derivative convolved, up to sign, with $\mathbf1_{[0,1]}^{*3}$.
Its convolution factor has $L^1$ norm one. Thus
$a_r\le\|h_r^{(3)}\|_{\TV}/16$ and
$b_r\le\|h_r^{(4)}\|_{\TV}/16$, giving
\[
\sum_{r=0}^{n-1}a_r\le\frac{\sqrt2}{16}\mathsf H_n,
\qquad
b_*\le\frac1{16}\sum_{r\ge0}
\min\{2^{-3},\sqrt6(r+1)^{-3/2}\}
\le\frac{3\sqrt[3]6}{32}.
\]
For the last sum, bound it by
$\int_0^\infty\min(2^{-3},\sqrt6 x^{-3/2})dx
=3\sqrt[3]6/2$.
Therefore \eqref{eq:maxW} and $e_n\le W_n$ give
\[
\Delta_n\le
\frac{16+\frac{\sqrt2}{8}\mathsf H_n}
     {1-\frac{3\sqrt[3]6}{16}}
<\frac{128}{5}+\frac3{10}\mathsf H_n,
\]
using $\sqrt2<3/2$ and $\sqrt[3]6<2$.
Finally the CDF of $\alpha-\beta$ is $1/16,-2/16,1/16$ on the unit
intervals starting at $-2,-1,0$. Equation~\eqref{eq:increments} gives
$\Delta_1=(1+2+1)/32=1/8$.
\end{proof}

\section{A conjectured universal logarithmic coefficient}\label{sec:coefficient}
The classification determines growth order, not the leading constant.
The increment identity suggests a stronger law. Write
$\delta\kappa=\int x^3D(dx)$ and suppose the common variance is $v>0$.
Let $\Phi$ and $\phi$ be the standard normal CDF and density, and set
\[
\Psi_n(x)=\Phi\!\left(\frac{x}{\sqrt{nv}}\right),\qquad
c(\alpha,\beta)=\frac{|\delta\kappa|}{3v\sqrt{2\pi e}}.
\]
\begin{conjecture}[Sharp logarithmic coefficient]\label{conj:sharp}
For finite pure dyadic action laws with common entropy $h$, common
information variance $v>0$, and $\kappa_P\ne\kappa_Q$,
\begin{equation}\label{eq:conjecture}
C_n^{\mathrm{causal}}(P,Q)
=nh+\frac{|\kappa_P-\kappa_Q|}{3v\sqrt{2\pi e}}\ln n+O(1).
\end{equation}
The remainder may depend on the pair.
\end{conjecture}
The coefficient is in bits per unit of $\ln n$. If the logarithm in
\eqref{eq:conjecture} is written as $\log_2 n$, the displayed coefficient
must be multiplied by $\ln 2$.

\paragraph{Formal origin of the constant.}
Since the first three moment cancellations include orders $0,1,2$,
Taylor expansion against a smooth profile gives
\[
D*\Psi_n=-\frac{\delta\kappa}{6}\Psi_n'''
             +\text{higher-order terms}.
\]
Moreover,
\[
\Psi_n'''(x)=\frac{1}{(nv)^{3/2}}
 \left(\frac{x^2}{nv}-1\right)\phi\!\left(\frac{x}{\sqrt{nv}}\right),
\qquad
\|\Psi_n'''\|_1=\frac{4\phi(1)}{nv}.
\]
For the last equality, integrate $(1-z^2)\phi(z)$ on $[-1,1]$,
using $(z\phi(z))'=(1-z^2)\phi(z)$, and use the zero integral of
$(z^2-1)\phi(z)$ on $\mathbb R$. Substitution into the exact increment
identity predicts $e_{n+1}-e_n\sim c(\alpha,\beta)/n$.
The missing step is control of the actual nonlinear envelope at the
scale of this differentiated $L^1$ expression. A central limit theorem
for the undifferentiated CDF alone would not justify that substitution.

\begin{proposition}[A sufficient remainder estimate]\label{prop:criterion}
Let $\alpha,\beta$ be finite centered laws of common variance $v>0$,
with $\delta\kappa\ne0$. Define
\[
\varepsilon_n=\left\|D*G_n+\frac{\delta\kappa}{6}\Psi_n'''\right\|_1,
\qquad n\ge1.
\]
Then
\begin{equation}\label{eq:conditional}
\left|(e_{n+1}-e_n)-\frac{c(\alpha,\beta)}n\right|
\le\frac{\varepsilon_n}{2}.
\end{equation}
If $\varepsilon_n=o(n^{-1})$, then
$e_n=c(\alpha,\beta)\ln n+o(\ln n)$.
If $\sum_{n\ge1}\varepsilon_n<\infty$, then
$e_n-c(\alpha,\beta)\ln n$ converges to a finite limit.
\end{proposition}
\begin{proof}
The reverse triangle inequality, \eqref{eq:increments}, and the computed
norm of $\Psi_n'''$ give \eqref{eq:conditional}. Sum over $n$.
An $o(1/n)$ sequence has partial sums $o(\ln n)$; absolute summability
gives convergence after subtracting the harmonic sum. Finally
$\mathsf H_n-\ln n$ converges.
\end{proof}
The proposition is conditional: neither of its remainder hypotheses is
proved here for the full class in Conjecture~\ref{conj:sharp}.
The summable estimate is stronger than what is needed for an $O(1)$
remainder, but it specifies one concrete sufficient analytic target.

\paragraph{Consistency checks.}
Several transformations constrain a proposed coefficient. Reflection
changes the sign of $\delta\kappa$ but not the growth cost. Dilation of
both centered increment laws by $a>0$ gives exactly
$G_n^{(a)}(x)=G_n(x/a)$ and $e_n^{(a)}=ae_n$; the proposed coefficient
has the same scaling since $v\mapsto a^2v$ and
$\delta\kappa\mapsto a^3\delta\kappa$.
Adding the same independent centered noise to each increment leaves
$\delta\kappa$ unchanged and raises $v$. The formula therefore predicts
a smaller logarithmic coefficient. This last prediction is a consequence
of the conjecture, not an additional theorem.

\section{Numerical evidence and exact finite certificates}\label{sec:numerics}
We evaluate the CDF recursion deterministically on its complete integer
grid; there is no sampling, fitted coefficient, or tail truncation.
Let $\alpha,\beta$ be the centered laws in Section~\ref{sec:explicit}, and
put
\[
\gamma=\tfrac18\delta_{-2}+\tfrac34\delta_0+\tfrac18\delta_2,
\qquad \check\alpha=\Law(-X),\quad X\sim\alpha.
\]
In Table~\ref{tab:coefficient}, $2\alpha$ denotes the pushforward under
$x\mapsto2x$, not multiplication of total mass. All listed laws are
centered with dyadic rational weights. Each pair therefore admits an
exact realization as centered information laws of finite pure dyadic
action laws: if all weights have denominator dividing $2^s$ and the
support lies in $[-K,K]\cap\mathbb Z$, choose $h=s+K$ and replace mass
$p_k$ at $k$ by $p_k2^{h+k}$ atoms of probability $2^{-(h+k)}$.
These are integer multiplicities, the total mass is one, and the entropy
is $h$.

\begin{table}[H]
\centering\small
\setlength{\tabcolsep}{4pt}
\begin{tabular}{@{}lrrrrr@{}}
\toprule
Increment laws & $v$ & $|\delta\kappa|$ & $c$ & $n(e_n-e_{n-1})$ & $10^4(r_n-1)$\\
\midrule
$\alpha,\beta$ & 1 & 0.375 & 0.03024634 & 0.03025164 & 1.752 \\
$(\alpha+\beta)/2,\beta$ & 1 & 0.1875 & 0.01512317 & 0.01512696 & 2.506 \\
$\alpha,\gamma$ & 1 & 0.375 & 0.03024634 & 0.03025556 & 3.047 \\
$\alpha,\check\alpha$ & 1 & 0.75 & 0.06049268 & 0.06049348 & 0.132 \\
$\alpha*\beta,\beta*\beta$ & 2 & 0.375 & 0.01512317 & 0.01512514 & 1.305 \\
$(\alpha+3\beta)/4,(3\alpha+\beta)/4$ & 1 & 0.1875 & 0.01512317 & 0.01512413 & 0.632 \\
$2\alpha,2\beta$ & 4 & 3 & 0.06049268 & 0.06050328 & 1.752 \\
\bottomrule
\end{tabular}
\caption{Deterministic floating-point recursion at $n=10\,000$.
Here $c=|\delta\kappa|/(3v\sqrt{2\pi e})$ is specified in advance and
$r_n=n(e_n-e_{n-1})/c$. The last row is an exact scaling consistency
check, not an independent family. Finite agreement does not establish
the asymptotic formula.}\label{tab:coefficient}
\end{table}

\begin{figure}[H]
\centering
\includegraphics[width=\textwidth]{coefficient_evidence.pdf}
\caption{Left: normalized entropy increments approach the predicted
coefficient over the displayed range. The cubic and dilation curves
coincide by exact scaling. Right: centered excess for three selected
families. Apparent stabilization over a finite range does not prove a
bounded remainder. Full parameters and numerical rows are supplied in
\texttt{numerical\_results.json} and \texttt{numerical\_results.csv}.}
\label{fig:evidence}
\end{figure}

For each of the seven pairs, integer arithmetic verifies normalization,
nonnegative profile masses, integer dyadic seed multiplicities, and the
entropy-increment identity through horizon $128$. Floating-point values
at that horizon are compared with the exact rational result.
For the cubic pair, a second run at $n=10\,000$ using the platform's
extended-precision \texttt{longdouble} changes $n(e_n-e_{n-1})$ by
approximately $2.1\times10^{-11}$ and $e_n$ by approximately
$1.1\times10^{-11}$. These comparisons check arithmetic sensitivity;
they do not bound the asymptotic remainder.

A separate exact program constructs the full dyadic seed partitions
through horizon two for the explicit action pair, and verifies the
conditional product law of every decoder branch. Its seed sizes are
$84$ and $6984$. At longer horizons it keeps atom multiplicities instead
of expanding all seed states. In particular, it verifies
\[
\Delta_1=\frac18,\qquad \Delta_2=\frac7{32},
\]
and the inequalities of Theorem~\ref{thm:explicit} at each of the $128$
checked horizons. The all-horizon claims follow from the proofs above,
not from these finite certificates.

\section{Scope and remaining questions}
The analytic envelope classification and the exact pure dyadic causal
classification are proved separately. Outside the pure dyadic class,
the present argument supplies necessary moment conditions for bounded
causal excess but no general matching seed construction. Even when the
analytic envelope has bounded mean, realization may introduce additional
cost. This realization question is distinct from the sharp-coefficient
question, which already arises within the class where attainment is exact.

The logarithmic regime is critical for the method: cancellation through
order two produces an $r^{-1}$ smoothing cost, while one further moment
cancellation makes that cost summable. The Fourier argument proves that
a genuine third-moment defect cannot be removed by a different seed
construction in the dyadic class. Higher moments can change the bounded
limit, but cannot create another divergent order in that class.

The numerical study tests a specific coefficient and its transformation
rules. A proof would require finer control of the nonsmooth nonlinear
recursion than the integrated contraction used for growth orders.
Proposition~\ref{prop:criterion} states one such target explicitly.
No conclusion about a common infinite source is implied: a different
finite seed is allowed at each horizon, and the expanded decoders can
have exponential size.

\appendix
\section{Constructing and checking a finite decoder}\label{app:decoder}
For a pure dyadic information profile $F_n$, let
\[
N_{n,k}=2^k\bigl(F_n(k)-F_n(k-1)\bigr).
\]
Create $N_{n,k}$ labeled seed atoms of mass $2^{-k}$ and sort them from
largest to smallest. For action $a$, the target bins are labeled pairs
$(y,z)$, with capacities $a(y)R_{n-1}(z)$. Fill these bins with the sorted
source atoms by the packing argument of Lemma~\ref{lem:optimal}. A seed
atom's bin label determines both the output $y$ and the next state $z$.
The exact identity
\[
\Prb\{Y=y,Z=z\mid A=a\}=a(y)R_{n-1}(z)
\]
proves that the next state has the required law conditional on the
observed output. Iterating gives the history-indexed simulator.

For the explicit pair, the certificate stores every source atom's
information length and its target-bin index for both actions at horizons
one and two. If the previous seed has $M$ atoms and the target index is
$t$, the output index and next-state index are respectively
$\lfloor t/M\rfloor$ and $t\bmod M$. All partition checks use integer
masses on a common dyadic grid.

For fixed finite dyadic action laws, the information support has width
$O(n)$. Thus the profile can be updated in $O(n)$ arithmetic operations
per horizon, using $O(n)$ stored profile values, without expanding the
seed. These are arithmetic-operation counts: exact numerator bit lengths
also grow with $n$. No polynomial bound on the expanded decoder size is
asserted.

\section{Reproduction}
The accompanying package includes the manuscript source, figure, exact
finite certificates, complete numerical tables, and two Python programs.
From the package directory, run
\begin{verbatim}
python3 experiments.py --n 10000 --exact-n 128
python3 logarithmic_seed.py --exact-n 128 --float-n 0
pdflatex -interaction=nonstopmode -halt-on-error moment_thresholds.tex
pdflatex -interaction=nonstopmode -halt-on-error moment_thresholds.tex
\end{verbatim}
The first program uses NumPy and Matplotlib; the second uses only the
Python standard library when \texttt{--float-n 0} is selected. The JSON
report records Python and NumPy versions and the available extended
precision. All probability parameters in the exact computations are
rational numbers. The predicted coefficient is computed directly from
those parameters, without regression.

\begin{thebibliography}{9}\small\raggedright\setlength{\itemsep}{4pt}
\bibitem{CGV19}
F. Cicalese, L. Gargano, and U. Vaccaro.
Minimum-entropy couplings and their applications.
\emph{IEEE Transactions on Information Theory}, 65:3436--3451, 2019.
\href{https://doi.org/10.1109/TIT.2019.2894519}{doi:10.1109/TIT.2019.2894519}.

\bibitem{Li21}
C. T. Li.
Efficient approximate minimum entropy coupling of multiple probability
distributions.
\emph{IEEE Transactions on Information Theory}, 67(8):5259--5268, 2021.
\href{https://doi.org/10.1109/TIT.2021.3076986}{doi:10.1109/TIT.2021.3076986}.

\bibitem{SY23}
Y. Y. Shkel and A. K. Yadav.
Information spectrum converse for minimum entropy couplings and functional
representations.
\emph{IEEE International Symposium on Information Theory}, 2023.
\href{https://arxiv.org/abs/2305.05745}{arXiv:2305.05745}.

\bibitem{CKQGK23}
S. Compton, D. Katz, B. Qi, K. Greenewald, and M. Kocaoglu.
Minimum-entropy coupling approximation guarantees beyond the majorization
barrier.
\emph{Proceedings of Machine Learning Research}, 206:10445--10469, 2023.
\href{https://proceedings.mlr.press/v206/compton23a.html}{PMLR paper and supplementary material}.

\bibitem{Compton26}
S. Compton.
Efficient $\varepsilon$-approximate minimum-entropy couplings.
\href{https://arxiv.org/abs/2509.19598v3}{arXiv:2509.19598v3},
18 June 2026.

\bibitem{Ma25}
Y.-J. Ma, F. Wang, and X.-Y. Wu.
Efficient approximate Minimum-R\'enyi entropy couplings.
\emph{Discrete and Continuous Dynamical Systems--S}, 18(10):2924--2936, 2025.
\href{https://doi.org/10.3934/dcdss.2025098}{doi:10.3934/dcdss.2025098}.

\bibitem{YS26}
A. K. Yadav and Y. Y. Shkel.
Geometry of R\'enyi entropy on the majorization lattice.
\href{https://arxiv.org/abs/2605.09655v2}{arXiv:2605.09655v2},
22 May 2026.

\bibitem{Guo24}
X. Guo.
On the rate of convergence of the martingale central limit theorem in
Wasserstein distances.
\href{https://arxiv.org/abs/2407.16980}{arXiv:2407.16980}, 2024.

\bibitem{FPSS19}
X. Fang, S. Peng, Q.-M. Shao, and Y. Song.
Limit theorems with rate of convergence under sublinear expectations.
\href{https://arxiv.org/abs/1711.10649v2}{arXiv:1711.10649v2}, 2018.
\end{thebibliography}
\end{document}
